\documentclass[11pt,a4paper]{article}

\usepackage[margin=25mm]{geometry}
\usepackage[T1]{fontenc}
\usepackage{lmodern}
\usepackage{amsmath,amssymb,amsthm}

\theoremstyle{definition}
\newtheorem{definition}{Definition}
\newtheorem{problem}{Problem}

\newcommand{\A}{A}
\newcommand{\klarner}{\lambda}

\setlength{\parindent}{0pt}
\setlength{\parskip}{6pt}

\begin{document}

\begin{center}
  {\Large\bfseries Klarner's Constant}
\end{center}
\section{Definitions}

\begin{definition}[Polyomino]
A \emph{polyomino} is a finite set $P\subset \mathbb Z^2$ of unit square cells that is
connected by edge adjacency.  Its \emph{area} (or size) is $|P|$; a polyomino of
area $n$ is called an $n$-omino.
\end{definition}

\begin{definition}[Fixed and free polyominoes]
Two polyominoes are \emph{fixed-equivalent} if they differ by a translation.  Let
$\A(n)$ be the number of fixed $n$-ominoes.  
\end{definition}

If rotations and reflections are also
identified, one obtains \emph{free polyominoes}.  Since every free class contains at
most eight fixed classes, fixed and free polyominoes have the same exponential growth
constant.  We focus on $\A(n)$.

\begin{definition}[Klarner's constant]
The \emph{Klarner constant} is
\begin{equation}\label{eq:klarner}
 \klarner
 =\lim_{n\to\infty}\A(n)^{1/n}
 =\lim_{n\to\infty}\frac{\A(n+1)}{\A(n)}.
\end{equation}
\end{definition}


\section{Current best result}
The exact value of $\klarner$ is unknown. The current rigorous interval is
\begin{equation}
  4.00253<\klarner\leq 4.5252.
\end{equation}
A stronger upper bound, $\klarner\leq 4.5238$, has also been reported.

\section{Problem formulation}
Improve the lower bound $4.00253$ or the reported upper bound $4.5238$.

\end{document}
